% Abstract in English here (1 page)

Methane is the second most abundant \acrlong{ghg} in the atmosphere (1935~ppb) after \ch{CO2} (423~ppm) and is responsible for around 30\% (\gls{gross}) of anthropogenic global temperature rise. Because of its high \acrlong{gwp}, especially on a short time period, methane is an excellent target for mitigation. Right now, methane is still accumulating globally in the atmosphere through the emission of mainly three sectors: energy, waste and agriculture. While the emission from the energy and the waste sector could be reduced with established technologies easily, the agriculture sector is classified as "hard-to-abate". This bachelor thesis investigates one possible mitigation strategy to solve this: the "immunization of cattle". For this purpose, a more holistic approach was chosen, combining a biochemistry and a geography part. 

In the biochemistry part, I performed a microbiological investigation of \textit{Methanobrevibacter} as the dominant rumen archaea, which are the biological entities that generate methane in the digestive tract of ruminants. The goal was to achieve DNA transfer in \textit{Methanobrevibacter}, which has not been demonstrated before. Thereby, the focus was primarily on two species (\acrlong{thau} and \acrlong{bovi}). It was shown that the species can be cultivated in pure culture in liquid and solid media and are sensitive to puromycin and \acrshort{dmso}. Furthermore, promising results were achieved to transform these species for the first time via interdomain conjugation with \textit{Escherichia coli}, but verification via whole-genome sequencing is still pending. In the future, the established genetic tool can help to identify a specific and potent vaccine that targets only the anaerobic, methanogenic archaea in the rumen of cattle. The long-time goal is to shift the microbiome towards more non-methanogenic microbes, such as acetogenic bacteria, so that the digestive efficiency of cattle is not compromised.

In the geographical part, I tried to evaluate the potential role of the investigated mitigation strategy to slow down climate change. For this purpose, the changed methane emission from enteric fermentation processes in cattle were calculated for different success scenarios of microbiome shift for three countries (Germany, New Zealand, USA) and its impact on national \ch{CH4} emission. The calculation is based on the 2006 published guidelines by the \acrshort{ipcc}. The geodata analyses showed for a shift of at least 30\% a relevant impact to achieve international emission targets. Given the current state of vaccine development, this method cannot play a significant role in reducing emissions over the next decade. However, in the long term, the technology could offer a solution that allows for a sustainable cattle population size that is socially and politically acceptable, even though a reduction will likely still be necessary.

Future research should further develop genetic tools to intensify the investigation of methanogenic archaea. Additionally, it should be assessed whether immunization of cattle could affect emissions beyond enteric methane, either amplifying or offsetting the overall emission reduction potential.

